\documentclass[10pt]{amsart}
\usepackage{amsmath,amssymb,amsthm}
\setlength{\textwidth}{\paperwidth}\addtolength{\textwidth}{-3.6cm}
\setlength{\oddsidemargin}{-0.74cm}\setlength{\evensidemargin}{-0.74cm}
\setlength{\headheight}{0pt}\setlength{\headsep}{0pt}
\setlength{\topmargin}{-0.84cm}
\setlength{\textheight}{\paperheight}\addtolength{\textheight}{-3.6cm}
\setlength{\footskip}{0.8cm}
\ifdefined\pdfpagewidth\pdfpagewidth=\paperwidth\pdfpageheight=\paperheight\fi
\pagestyle{plain}\raggedbottom
\setlength{\parskip}{0pt plus 0.5pt}
\widowpenalty=150 \clubpenalty=150
\makeatletter\def\section{\@startsection{section}{1}{\z@}{.6\baselineskip}{.25\baselineskip}{\normalfont\bfseries}}\makeatother
\newtheorem{prop}{Proposition}
\newtheorem{cor}[prop]{Corollary}
\newcommand{\st}[1]{\par\noindent\textit{#1}\enspace}

\title{What Section 6 gives at slope $\pm2$}
\date{}
\begin{document}\sloppy
\maketitle
\vspace{-2.5em}

\noindent\textbf{Outcome.} I could not derive $S^3_2(K)\not\cong S^3_{-2}(K)$ from
Section~6, even with nondegeneracy hypotheses. Section~6 does give an exact
formula for every drop, with the window ends as corollaries. The formula shows
precisely what is missing: the position of the class realising the spectral
invariant. I see no hypothesis of transversality type that controls this position,
and the data show that it is not controlled for free.

\section{Set-up}

$K\subset S^3$ satisfies (H1)--(H3) and $\Delta_K''(1)=0$. The latter is forced by a
cosmetic pair through the Casson--Walker invariant. Let
$U=(0,\frac12)_a\times(\mathbb R/\mathbb Z)_b$ be the pillowcase minus the edges
$a\in\{0,\frac12\}$, with $\omega=2\,da\wedge db$. By Proposition~1.1(b),
$\Gamma=\Gamma_K\subset U$. The surgery lines $L_N^\varepsilon=\{Na+b\equiv\varepsilon/2\}$
are arcs crossing $U$. Since $\Gamma$ misses the edges, all these arcs are
interchangeable as far as $\Gamma$ is concerned.

\begin{prop}[from Corollary 6.4; minimal position assumed]
$CI(Y_{\pm2},\varepsilon)\cong CI(Y_{\pm1})\cong CI(Y_0,w)$ as $\mathbb Z/8$-graded
complexes, up to grading shifts, for both $\varepsilon$. The $S^3$ terms vanish. In
particular there is no unfiltered obstruction: the problem lives entirely in the
Chern--Simons data. Moreover $\Gamma$ has total winding $0$ around $U$. By Herald's formula (\S3.1),
the algebraic count of $\Gamma\cap L_{1/n}$ is $2\lambda(Y_{1/n})=n\Delta''_K(1)=0$.
\end{prop}

\section{The exact drop formula}

Let $\lambda_{1/4}=2(a-\frac14)\,db$. This primitive of $\omega$ is invariant under
$\chi(a,b)=(\frac12-a,-b)$. Let $\Phi$ be its primitive along $\Gamma$, which is
$d\Phi=2(a-\frac14)db$ of \texttt{status.pdf}. On $L_N$ the restriction of
$\lambda_{1/4}$ is $d\bigl(-N(a-\frac14)^2\bigr)$. With the line primitive normalised to
vanish at its traceless point $a=\frac14$, the pillowcase level formula reads
\[
c(Y_N,x)=A_N(x)+\tfrac3{16}\operatorname{sgn}N-\tfrac N{16}-\tfrac18\iota_N(x),\qquad
A_N(x)=\Phi(x)+N\bigl(a(x)-\tfrac14\bigr)^2 .
\]
This is exact for torus knots and checked numerically for $4_1$. For $N=\pm2$ the
constant is $\pm\frac1{16}$: this is where the sixteenth comes from.

\begin{prop}[rotation identity]
Let $P=(a_P,b_P)$ be a common point of $L_2^\varepsilon$ and $L_{-2}^{\varepsilon'}$. Let
$L_t$, $t\in[-2,2]$, be the lines through $P$ of slope $-t$, and let $x_t\in\Gamma\cap L_t$
move continuously, with no birth or death. Then the level moves at the rate
$-(a_t-a_P)^2-a_P(\frac12-a_P)$. Hence
\[
c(Y_2,x_2)-c(Y_{-2},x_{-2})=g_P-\int_{-2}^{2}\bigl(a(x_t)-a_P\bigr)^2dt-\tfrac18\Delta\iota,
\]
where $g_P=\frac18$ for the same bundle ($P=(\frac14,\frac{\varepsilon-1}2)$) and
$g_P=\frac3{16}$ for mixed bundles ($a_P\in\{\frac18,\frac38\}$).
\end{prop}

\noindent\emph{Proof.} On $\{t(a-a_P)+b-b_P=0\}$ we have
$\dot b=-(a-a_P)-t\dot a$. Hence $\frac d{dt}A_t=2(a-\frac14)\dot b+(a-\frac14)^2+2t(a-\frac14)\dot a$.
The terms in $\dot a$ cancel, leaving $(a-\frac14)(2a_P-a-\frac14)$. Adding the
$-\frac1{16}$ of the constant gives the stated rate. The constant jumps by
$\frac38$ at $t=0$. At $x=P$ the integral vanishes and the gap is $g_P$. $\square$

\noindent This is Section~6's $\delta(x)$ written as a swept area. For a corner
pivot it is literally Theorem~6.3: the triad $(1,\frac10,2)$ has midpoint
$m_1=(\frac12,\frac12)$, a corner. \emph{Check:} for $T(2,3)$ the identity
the level formula gives $u=\frac3{32}$ and $w=\frac14$, as in the tables.
Here $\Gamma$ is the segment $b=\frac12-6a$ and $\Phi=-6(a-\frac14)^2$.
This is a weak check: the formula was calibrated on torus knots.

\begin{cor}[windows]
Assume continuity of spectral invariants along the family (standard for the
holonomy perturbations of \S5.3), and that the grading term $\Delta\iota$
vanishes along it, as in every example. Then every drop equals $g_P$ minus the integrated
squared distance of the \emph{spectral carrier} from $P$. Hence the same-bundle
drop lies in $(-\frac18,\frac18]$ and the mixed drop in $(-\frac38,\frac3{16}]$. The
upper ends are attained exactly when the carrier sits at the common point.
\end{cor}

\noindent These are the known windows: DLME's map $H$ and the map $B$ for the same
bundle, and the third map for mixed bundles. The equality cases are the
common-point cases already seen in the data.

\section{Why this does not prove the theorem}

A cosmetic pair needs every drop to vanish: bottom and top, both bundles. By the
identity this means that each carrier's integrated squared distance from its
pivot equals $g_P$. For the same bundle, $\int_{-2}^2(a_t-\frac14)^2dt=\frac18$:
an RMS distance of $\frac1{4\sqrt2}$ from traceless. Nothing in Section~6 excludes
this:
\begin{itemize}
\item $\Gamma$ is an arbitrary exact, $\chi$-invariant multicurve of winding $0$ in $U$;
\item the cone isomorphisms only transport generators along such rotations;
\item transversality, minimal position and distinct levels say nothing about
where the carrier sits.
\end{itemize}
The data show that it is not controlled for free. For positive torus knots the
trivial-bundle drop $u$ is negative, so that carrier's deficit exceeds $\frac18$.

A hypothesis bounding the carrier's distance from $P$ does suffice. For example,
if $\Gamma\subset\{|a-\frac14|<\frac{\sqrt2}8\}$, every same-bundle drop is positive.
But such a bound is a bound on the Hamiltonian along the carrier, i.e.\ an energy
hypothesis in disguise. This is the input the no-go theorem said was missing:
\emph{where the surviving class sits.}

\section{Where a lower bound can come from}

Lower bounds do not come from rotations through the horizontal, i.e.\ through
$Y_{\pm1}$, $Y_0$ and the cones of Section~6. They can only come from the other
route, through the traceless circle $\{a=\frac14\}$: the bundle-exchanging triangle
with the theory of $(S^3,K)$, which is not a Farey triad. There a class that is not
absorbed has mixed drop greater than $\frac18$ (the kernel criterion of
\texttt{triangle-1}, conditional on its hypotheses). Absorption
is the only failure mode, and Section~6 says nothing about it. Zero-energy
absorption needs a common traceless point with $\lambda\mapsto\pm1$. For
$\Delta_K=1$ these are not binary dihedral and occur only degenerately, as for
$P(-3,5,7)$.

\section{Calibration at $\pm1$, where the answer is known}

$L_1$ and $L_{-1}$ meet only at the corners $(0,0)$ and $(\frac12,\frac12)$. Section~6
joins $Y_{\pm1}$ through $(Y_0,w)$ by the triads $(0^w,\frac10,\pm1)$, which is
rotation about the corner $C=(\frac12,\frac12)$. The same computation gives
\[
\ell(Y_1)-\ell(Y_{-1})=\tfrac38-\operatorname{Area}(T)-\tfrac18\Delta\iota ,
\]
where $T$ is the region swept by the carrier, with vertex $C$. \emph{Check:} for the
right-handed trefoil the level formula gives $\frac1{80}-(-\frac{23}{112})=\frac{61}{280}$,
which is DLME's $\frac{59}{120}-\frac{23}{84}$; here $\Delta\iota=0$ and $\operatorname{Area}(T)=\frac{11}{70}$.
The window $(-\frac38,\frac38)$ contains $0$, so \emph{Section~6 does not prove the
$\pm1$ case either.}

DLME's lower bound $\ell(Y_1)-\ell(Y_{-1})\ge\frac18+\eta$ comes from their map
$g_1\colon Y_1\to Y_{-1}$: a one-parameter family with $b^+=0$ and level $-\frac18-E$, which is
an isomorphism because $I(S^3)=0$. Take a pair $x\mapsto y$ that $g_1$ counts, with
$x,y$ on one component of $\Gamma$. Stokes gives
\[
E=\tfrac14-\operatorname{Area}(T_{xy})-\tfrac18\Delta\iota .
\]
So positivity of $g_1$'s energy \emph{is} the area bound that the cones cannot see.
At $\pm2$ the map in $g_1$'s role is $B$, of level $+\frac18$. Positivity of its
energy gives only a drop greater than $-\frac18$.

\end{document}
